\chapter{Ada Language Techniques}
\label{ch:embedded}

A handful of Ada features come up again and again in embedded code: writing
literals clearly, constraining values to only what is legal, validating data that
arrives from outside, reinterpreting bits deliberately, attaching checkable
contracts, and handling (or forbidding) exceptions. This chapter collects the
language-level techniques; the next, \emph{Talking to Hardware}, applies them to
registers and memory. Together they underpin the HAL and the filesystem, so a few
terms are worth learning before the driver chapters that use them.

\section{Terminology}
\label{sec:terms}\index{terminology}

A few Ada words recur throughout the book and can trip up a reader coming from C.
They are defined once here.

\begin{description}[style=nextline]
\item[Access type / \texttt{'Access}.] Ada's typed pointer; \texttt{X'Access}
  yields a pointer to \texttt{X}. Strongly typed, with accessibility rules that
  guard against dangling references.
\item[Activation.] The moment a declared task begins executing: after its master's
  declarations, or, for a library-level task, as part of elaboration before
  \texttt{Main}.
\item[Aspect.] A property attached to a declaration with the word \texttt{with}:
  \texttt{with Volatile}, \texttt{with Address => \ldots}, \texttt{with Pre =>
  \ldots}. Comparable to a GCC \texttt{\_\_attribute\_\_}, but standardised and
  checked.
\item[Attribute.] A built-in query or operation written with an apostrophe:
  \texttt{X'Image}, \texttt{T'Last}, \texttt{Obj'Address}, \texttt{P'Access}. There
  is no C equivalent.
\item[Ceiling priority.] The priority at which a protected object's operations
  run (the priority-ceiling protocol); for an interrupt handler it is the
  interrupt's hardware level.
\item[Constraint check / \texttt{Constraint\_Error}.] The run-time check that a
  value is in range or valid; a failure raises \texttt{Constraint\_Error}. The
  checks can be switched off in a release build.
\item[Elaboration.] The once-only, ordered set-up that runs \emph{before}
  \texttt{Main}: evaluating initial values, executing the statements in package
  bodies, and activating library-level tasks. Like C's static initialisers, but
  ordered and able to run arbitrary code.
\item[Entry / barrier.] A protected (or task) operation guarded by a Boolean
  condition, the \emph{barrier}; a caller blocks until the barrier is open. This is
  how condition synchronisation is expressed.
\item[Generic.] A compile-time template (a package or subprogram parameterised
  by types or values) instantiated to produce a concrete unit, much like a C++
  template.
\item[Library unit / library-level.] A unit (package or subprogram) declared at
  the outermost level, \emph{not} nested inside a subprogram or block. A
  library-level object or task exists for the whole run of the program. C has only
  file scope and block scope; Ada nests arbitrarily, so ``library-level'' names
  the top, program-lifetime level, and the Jorvik profiles require tasks (and
  some other entities) to live there.
\item[Package (specification / body).] A named module. The \emph{specification}
  is the visible interface (roughly a C header) and the \emph{body} is the
  implementation (roughly the \texttt{.c}), but they are one logical unit and
  the body's internals are genuinely hidden by the compiler, not merely by
  convention.
\item[Pragma.] A directive to the compiler: \texttt{pragma Pack},
  \texttt{pragma Attach\_Handler}, \texttt{pragma Restrictions}. Like
  \texttt{\#pragma}, but part of the language with defined meanings.
\item[Profile (Ravenscar / Jorvik).] A named bundle of \texttt{Restrictions}
  defining an analysable tasking subset. Distinct from this project's runtime
  \emph{build} profiles, which choose \emph{which} such subset the runtime
  supports.
\item[Protected object.] Data plus operations with \emph{automatic} mutual
  exclusion: a monitor. Access is race-free without a manual lock; the compiler
  generates the locking.
\item[Rendezvous.] A synchronous call from one task to another task's
  \emph{entry}: the two meet, exchange parameters, the accept body runs, then both
  continue (the \texttt{full} profile).
\item[Representation clause.] A clause that fixes a type's machine layout (bit
  positions, size, byte order): \texttt{for T use record \ldots\ end record}. It is
  how a record is mapped onto a hardware register.
\item[Secondary stack.] A runtime area used to return a result whose size the
  caller does not know in advance (e.g.\ a function returning \texttt{String}). C
  has no equivalent; the leanest profile omits it.
\item[Task.] A thread of control defined in the language. It \emph{activates}
  (starts) automatically. There is no explicit start call.
\item[Type vs.\ subtype.] A \emph{type} is a distinct set of values and
  operations; two distinct types never mix without an explicit conversion. A
  \emph{subtype} is the same type plus a constraint (a range or predicate),
  fully compatible with its base, just narrower.
\item[Unconstrained type.] A type whose size is settled only when an object is
  created: \texttt{String}, \texttt{array (Positive range <>)}, a record with
  defaulted discriminants.
\end{description}

%  The lexical basics (numeric bases, underscores, and case insignificance)
%  are introduced early, with the first program, in \S\ref{sec:bases} (the Anatomy
%  chapter), so the register code in the intervening chapters reads correctly.
%  This chapter picks up where that leaves off: the TYPES those numbers take.

\section{Subtypes and modular types}
\label{sec:subtypes}\index{subtype}

The other way the type carries intent is by constraining which values are legal
at all. Two facilities earn their keep constantly in embedded code: subtypes that
narrow a value to a valid range or set, and modular types that give wrap-around
arithmetic and bit operations.

\subsection{Range subtypes: only the legal values}

A subtype attaches a constraint to an existing type. The everyday embedded use is
to make an out-of-range value impossible:

\begin{lstlisting}
   subtype Channel is Natural range 0 .. 7;     -- a peripheral with 8 channels
   subtype Percent is Natural range 0 .. 100;   -- a duty cycle, say
\end{lstlisting}

A value known at compile time that violates the range is a compile error; a value
computed at run time is caught by a constraint check (wherever checks are
enabled). Either way ``channel 9'' never reaches the hardware, and the range is
documented in the type rather than in a comment.

\subsection{Static-predicate subtypes: a non-contiguous valid set}

When the legal values are not one contiguous range, a \emph{static predicate}
states the exact set. The HAL's pin type is the real example: of the 0\,..\,48 pad
numbers, some are absent on the package and some are bonded to the in-package SPI
flash and octal PSRAM. Driving any of those hangs the chip:

\begin{lstlisting}
   type Pad_Number is range -1 .. 48;            -- -1 is the "no pin" sentinel

   subtype Pin_Id is Pad_Number range 0 .. 48 with
     Static_Predicate => Pin_Id in 0 .. 21 | 38 .. 48;   -- 22..37 reserved
\end{lstlisting}

Because the predicate is \emph{static}, naming a reserved pad as a constant of
type \texttt{Pin\_Id} is rejected at compile time (\emph{static expression fails
static predicate check}); a pin computed at run time is verified by the predicate
under \texttt{-gnata}. Every driver names its pins with \texttt{Pin\_Id}, so a
flash/PSRAM pad cannot be wired by accident.

\cnote{a pin is an \texttt{int} or a \texttt{\#define}, and wiring a flash pad is a
run-time hang you debug with an oscilloscope. \texttt{Pin\_Id} turns a bad
constant into a \emph{compile} error and a bad computed value into a checked one.}

\subsection{Modular types: wrap-around and bits}

A modular type (\texttt{mod 2**n}) holds the values $0\,..\,2^n-1$, and its
arithmetic \emph{wraps} instead of overflowing. That is exactly the behaviour you
want for a register word, a free-running counter, or a checksum; modular types
also support the logical operators \texttt{and}, \texttt{or}, \texttt{xor} and
\texttt{not}:

\begin{lstlisting}
   type Byte is mod 2 ** 8 with Size => 8;       -- 0 .. 255, wraps at 256

   Sum : Byte := 0;
   ...
   Sum   := Sum + B;                       -- no overflow error; wraps (a checksum)
   Flags := Flags or  2#0000_0100#;        -- set a bit
   Flags := Flags and not 2#0000_0100#;    -- clear a bit
\end{lstlisting}

The predefined \texttt{Interfaces.Unsigned\_8/16/32/64} are the usual choice for
register-sized words: they are modular and add \texttt{Shift\_Left} /
\texttt{Shift\_Right} / \texttt{Rotate\_*}. A plain signed integer would instead
raise \texttt{Constraint\_Error} on overflow and offers no bit operators. So the
rule of thumb is: registers, masks and counters are \emph{modular}, while
quantities that must never silently wrap (an index, a length) stay
\emph{range-constrained}.

\subsection{And beyond}

Two more worth keeping in mind:
\begin{itemize}
  \item \textbf{Fixed-point types} (\texttt{type Q8 is delta 1.0 / 256.0 range
        0.0 .. 1.0}) give fractional values with exact, integer-backed
        arithmetic: a duty cycle or a scaled sensor reading without dragging in
        floating point.
  \item \textbf{Derived types} make otherwise-identical numbers
        \emph{incompatible}: with \texttt{type Hz is new Natural;} and
        \texttt{type Millis is new Natural;}, a frequency can never be passed
        where a delay was meant. The mix-up is a compile error, not a bug.
\end{itemize}

\section{Arrays and arbitrary index ranges}
\label{sec:arrays}\index{arrays}\index{array index range}

An Ada array type fixes not just its element type but its \emph{index range}, and
that range is yours to choose: any discrete range at all. The lower bound need
not be 0 or 1; it can be any value (even negative), and the index can be an
enumeration type rather than an integer. The range is part of the type:

\begin{lstlisting}
type Sine_Table is array (0  .. 255) of Integer;     -- 0-based lookup table
type CAN_Data   is array (1  .. 8)   of Unsigned_8;  -- 1-based, eight bytes
type Window     is array (-3 .. 3)   of Float;       -- centred on zero
type Channel    is (Red, Green, Blue);
type Mix        is array (Channel)   of Unsigned_8;  -- indexed by an enum
\end{lstlisting}

The point is not novelty: it lets the index \emph{match the thing it models}, so
the code reads like the datasheet or the protocol, with no mental ``subtract
one'' and none of the off-by-one bugs that come with it.

\textbf{Iterating without knowing the bounds.} Every array value carries its own
bounds, available as attributes: \texttt{A'First}, \texttt{A'Last},
\texttt{A'Length} and \texttt{A'Range} (which is \texttt{A'First\,..\,A'Last}). The
idiomatic loop ranges over \texttt{'Range}, so it is correct whatever the bounds
are, and stays correct if you later change them:

\begin{lstlisting}
Total : Natural := 0;
for I in Buf'Range loop          -- I takes each actual index value in turn
   Total := Total + Natural (Buf (I));
end loop;
\end{lstlisting}

The same code works for a \texttt{0\,..\,255} buffer and a \texttt{1\,..\,8} one
because the loop names \texttt{Buf'Range}; there is no literal bound to get
wrong. (When you do not need the index itself, \texttt{for E of Buf loop}
iterates the elements directly.)

\textbf{A worked example: J1939 CAN payloads, indexed from 1.} A Classic CAN
frame carries up to eight data bytes. SAE J1939 (the protocol on commercial
vehicles) numbers those bytes \emph{1 through 8} in its specifications
(``byte 1,'' ``bytes 4--5,'' and so on). Declaring the payload with a
\texttt{1\,..\,8} index makes the Ada line up with the document exactly:

\begin{lstlisting}
type J1939_Payload is array (1 .. 8) of Interfaces.Unsigned_8;

--  PGN 61444 (EEC1): engine speed is SPN 190, in BYTES 4-5,
--  little-endian, scaled 0.125 rpm/bit.  The byte numbers here are
--  taken verbatim from the J1939 spec -- no off-by-one translation.
function Engine_RPM (D : J1939_Payload) return Float is
   Raw : constant Interfaces.Unsigned_16 :=
            Interfaces.Unsigned_16 (D (4)) +              -- byte 4: low
            Interfaces.Unsigned_16 (D (5)) * 256;         -- byte 5: high
begin
   return Float (Raw) * 0.125;
end Engine_RPM;
\end{lstlisting}

\texttt{D (4)} \emph{is} ``byte 4'' (the same number a J1939 reference uses),
so the code can be checked against the standard line by line, and the loop idiom
still applies: \texttt{for I in D'Range loop} walks bytes \texttt{1\,..\,8}.

\cnote{In C the eight bytes are \texttt{data[0..7]}, so every J1939 byte number
from the spec has to be mentally decremented as you read or write the code (spec
``byte 4'' becomes \texttt{data[3]}). An arbitrary lower bound removes that whole
translation step.}

\section{The \texttt{others} choice: filling in ``everything else''}
\label{sec:others}\index{others}

\texttt{others} is a catch-all \emph{choice}: it stands for ``every case not
named so far,'' and it must always come \emph{last}. Ada uses it wherever you
enumerate possibilities: array and record aggregates, \texttt{case} statements,
and exception handlers. It is everywhere in embedded Ada (it fills the reserved
fields of the register writes in Chapter~\ref{ch:hardware}), so it is worth a
close look.

\textbf{In array aggregates: fill the rest.} An aggregate must give a value for
\emph{every} element; \texttt{others} supplies the ones you did not name.

\begin{lstlisting}
Zeros  : Buffer := (others => 0);                            -- every element 0
Header : Buffer := (0 => 16#AA#, 1 => 16#55#, others => 0);  -- first two, rest 0
\end{lstlisting}

\textbf{In record aggregates, and the \texttt{<>} ``box''.} A record aggregate
must also set every component. Writing \texttt{others => <>} means ``give each
remaining component its \emph{default} value'': the default from the type's
own declaration. This is the central idiom for writing a hardware register: set
the one or two fields you mean, and let every reserved or don't-care field take
its declared default.

\begin{lstlisting}
--  A register record whose fields mostly default to 0 / False:
R.TX_CONF := (TX_START => True, others => <>);   -- set one bit, default the rest
\end{lstlisting}

Use \texttt{others => <>} rather than \texttt{others => 0} for a record whose
remaining components are of \emph{different} types: each takes its own default, so
the types need not agree. \texttt{others => 0} only works when everything it
covers is the same type (e.g.\ an all-zero byte buffer).

\textbf{In \texttt{case}: cover the remaining values.} A \texttt{case} must be
exhaustive; \texttt{when others} handles whatever you did not list.

\begin{lstlisting}
case Mode is
   when Standard => ...;
   when others   => ...;     -- every other value of the type
end case;
\end{lstlisting}

\textbf{In exception handlers: catch the rest.} \texttt{when others} catches
any exception not named earlier in the handler.

\begin{lstlisting}
begin
   ...
exception
   when Constraint_Error => ...;   -- a specific one
   when others           => ...;   -- anything else
end;
\end{lstlisting}

Two properties tie these uses together: \texttt{others} is always the \emph{last}
choice, and it covers \emph{exactly} the cases not already named, so adding a
component, value or exception elsewhere silently changes what ``everything else''
means. For a register write that is precisely what you want (a newly-defined
reserved bit still defaults). For a \texttt{case} over an enumeration it is often
the opposite: omitting \texttt{others} is safer because then \emph{adding} a value
to the type forces the compiler to make you handle it.

\section{Fixed-point arithmetic}
\label{sec:fixed}\index{fixed-point types}

Microcontroller code is full of fractional quantities (a duty cycle, a gain, a
scaled sensor reading) that do not justify (or cannot afford) floating point.
A \emph{fixed-point} type gives fractions with exact, integer-backed arithmetic:
you state the resolution with \texttt{delta}, and the compiler represents values
as integers scaled by that \texttt{'Small}.

\begin{lstlisting}
   type Duty is delta 1.0 / 256.0 range 0.0 .. 1.0;   -- 8-bit fraction, 0 .. 1
   D : Duty := 0.75;                 -- stored as the integer 192 (= 0.75 * 256)

   type Volts is delta 0.001 range 0.0 .. 3.3;        -- millivolt resolution
\end{lstlisting}

Arithmetic on \texttt{Duty} is integer arithmetic underneath (no FPU, no
rounding surprises), yet the source reads in real-world units. The PWM and ADC
drivers scale this way: a duty cycle or a voltage is a fixed-point value the
driver converts to the register's raw counts.

\section{Trusting external data: the \texttt{'Valid} attribute}
\label{sec:valid}\index{Valid attribute}

Anything that enters from outside the program (a value read from a register, a
field unpacked from an SD sector, a message off the wire) might not be a legal
value of its type. Reading it into a constrained type and then using it is
\emph{erroneous} if the bits are out of range. \texttt{X'Valid} asks ``is the
object actually a valid value of its (sub)type?'' and answers with a plain
Boolean (no exception, safe to ask even of garbage):

\begin{lstlisting}
   Raw : Parity with Address => Reg'Address, Volatile;  -- enum from hardware
begin
   if Raw'Valid then
      Use_Parity (Raw);             -- known to be None/Even/Odd
   else
      Handle_Corrupt;               -- some other bit pattern
   end if;
\end{lstlisting}

This is the safe way to consume an enumeration with reserved codes, a range field
that hardware might leave uninitialised, or a structure read from possibly-corrupt
storage: validate first, trust after.

\section{\texttt{Unchecked\_Conversion}: deliberate reinterpretation}
\label{sec:uncheck}\index{Unchecked\_Conversion}

Occasionally you must view the same bits as two different types: serialise a
record to bytes, decode a fat pointer, read a float's bit pattern.
\texttt{Ada.Unchecked\_Conversion} is the explicit, greppable way to do it; unlike
a C cast it is a named instantiation, so every reinterpretation in the codebase is
searchable.

\begin{lstlisting}
   function To_Fat is new Ada.Unchecked_Conversion (Handler_Access, Fat_Ptr);
\end{lstlisting}

The runtime uses exactly this to inspect an \texttt{access protected procedure}
as a \texttt{(object, code)} pair when matching a registered interrupt handler.
The discipline: prefer a representation clause or an address overlay when the two
views have a defined relationship; keep \texttt{Unchecked\_Conversion} for the
genuine bit-reinterpretation cases, where its name documents the hazard.

\section{Access types: where you need a pointer, and where you don't}
\label{sec:access}\index{access types}

Ada has pointers (\emph{access types}), but a great many of the places a C
programmer reaches for one do not need it. The embedded habit is to lead with the
pointer-free option and use an access type only where the problem genuinely
demands it.

\subsection{Where you can avoid them}

\textbf{Parameters are already passed efficiently.} You do not pass
\texttt{\&thing} and dereference. A parameter mode (\texttt{in},
\texttt{out}, \texttt{in out}) says how it is used, and the compiler passes a
large object by reference automatically:

\begin{lstlisting}
   procedure Process (Frame : in out Packet);   -- no pointer; Frame is the object
\end{lstlisting}

\textbf{Whole arrays and records are values.} You can assign, return and pass
entire structures; there is no need for a pointer just to ``return a struct.''

\begin{lstlisting}
   R := S;                          -- copies a whole record
   function Identity return Matrix; -- returns a whole array by value
\end{lstlisting}

\textbf{Unconstrained types size themselves.} A \texttt{String} or an
\texttt{array (Positive range <>)} takes its bounds from its value, so a function
can return one without \texttt{malloc} and without a pointer (the secondary stack
carries it):

\begin{lstlisting}
   function Pin_Name (P : Pin_Id) return String;   -- no heap, no pointer
\end{lstlisting}

\textbf{A register is an address overlay, not a pointer.} As
\S\ref{ch:registers} showed, you place a typed object \emph{at} the register's
address; you never hold a \texttt{volatile uint32\_t *} and dereference it.

\textbf{A resource is a handle, not a raw pointer.} The limited-controlled handles
of Chapter~\ref{ch:encap} own a peripheral by scope, not by a pointer you must
remember to free.

\subsection{Where you genuinely need one}

\begin{itemize}
  \item \textbf{Dynamic or linked data built at run time:} a list, a tree, a
        variable-length buffer pool: \texttt{type Node\_Ref is access Node;} and
        \texttt{new Node}.
  \item \textbf{Callbacks and interrupt handlers:} an access-to-subprogram
        value lets you register code to be called later. This is the
        \texttt{On\_Edge'Access} from the GPIO counter (\S\ref{sec:tasks}) and the
        handler registration in the Interrupts chapter.
  \item \textbf{Hardware structures that reference a buffer:} a DMA descriptor
        holds the \emph{address} of the data to move, so the buffer is
        \texttt{aliased} and the descriptor stores its \texttt{'Access}
        (\S\ref{sec:placement}).
  \item \textbf{Talking to C:} C interfaces are expressed in pointers, so
        \texttt{System.Address} and access types cross that boundary
        (\S\ref{sec:interfacing}).
\end{itemize}

\subsection{Why the pointer you do use is safer}

Even when you need one, an Ada access type is far less hazardous than a C pointer:

\begin{itemize}
  \item it is \textbf{named and typed:} a \texttt{Node\_Ref} is not a
        \texttt{Byte\_Ref} and neither is a \texttt{void *}; the compiler keeps
        them apart;
  \item there is \textbf{no pointer arithmetic:} you index an array, you do not
        increment a pointer past its end;
  \item \textbf{accessibility rules} reject \texttt{Local'Access} returned from the
        subprogram that owns \texttt{Local}, so a dangling reference is a compile
        error, not a crash;
  \item \texttt{\textbf{not null}} access subtypes exclude the null case at the
        type level;
  \item and \texttt{\textbf{for T'Storage\_Size use 0;}} (\S\ref{sec:memory})
        forbids \texttt{new} on an access type outright: a compile-time promise
        that this pointer type never touches a heap, which is how the
        \texttt{light-tasking} profile stays heap-less.
\end{itemize}

So the rule of thumb: prefer value semantics and address overlays; reach for an
access type for the genuinely dynamic structure, the callback, or the
hardware-referenced buffer, and even then it is typed, arithmetic-free, and, if
you wish, provably heap-free.

\section{Contracts for hardware-safe APIs}
\label{sec:contracts}\index{contracts}

Finally, a driver's \emph{interface} can carry its rules as checkable contracts,
so misuse is caught at the call site rather than producing a wrong register write.
Preconditions and postconditions state what must hold on entry and exit; type
invariants and predicates state what must always hold of a value. The AES driver,
for instance, accepts only a legal key length:

\begin{lstlisting}
   function Supported_Key (Key : Byte_Array) return Boolean is
     (Key'Length in 16 | 32);            -- AES-128 or AES-256 only

   procedure Set_Key (Key : Byte_Array)
     with Pre => Supported_Key (Key);    -- a 24-byte key is rejected at the call
\end{lstlisting}

The cost is tunable: \texttt{pragma Assertion\_Policy} (and the \texttt{-gnata}
switch) turns the run-time checks on for development and \emph{off} for a release
build, where the contracts remain as machine-checked documentation. Combined with
the range and predicate subtypes of \S\ref{sec:subtypes}, the effect is that a
large class of driver-misuse bugs (a bad pin, a wrong-width key, an
out-of-range channel) cannot reach the hardware.

\section{Exceptions and error handling}
\label{sec:exceptions}\index{exceptions}

An Ada \emph{exception} is a named error condition: \texttt{raise} signals it and
a handler catches it. The mechanics are familiar:

\begin{lstlisting}
   begin
      Read_Sector (N, Buf);
   exception
      when Storage_Error =>           -- out of memory / stack
         Recover;
      when E : others =>              -- anything else; E names it
         Log (Ada.Exceptions.Exception_Message (E));
   end;
\end{lstlisting}

A raise propagates up the call chain to the nearest matching handler, running
finalizers (Chapter~\ref{ch:encap}) on the way, which is why a controlled
handle is released even when an exception unwinds past it. If no handler matches,
it keeps propagating; if it escapes the whole program, the runtime's
\emph{last-chance handler} runs (below).

\textbf{Predefined and user-defined exceptions.} The language defines a handful
the runtime raises for you: \texttt{Constraint\_Error} (a failed range, index,
\texttt{null} or overflow check, what a bad \texttt{Pin\_Id} or an out-of-range
index produces), \texttt{Program\_Error} (an illegal runtime situation),
\texttt{Storage\_Error} (out of stack or heap), and \texttt{Tasking\_Error}. You
declare your own just as easily, and raise it with a message:

\begin{lstlisting}
   Not_Owned : exception;          -- a named error condition of your own
   ...
   if not Holding then
      raise Not_Owned with "use the bus without holding it";
   end if;
\end{lstlisting}

A handler can inspect the exception through \texttt{Ada.Exceptions}
(\texttt{Exception\_Name}, \texttt{Exception\_Message},
\texttt{Exception\_Information}), and a bare \texttt{raise;} inside a handler
\emph{re-raises} the current one, so you can clean up locally and still let an
outer scope deal with it.

\textbf{How this project uses them.} The ownership-checked HAL gateways raise
\texttt{Not\_Owned} / \texttt{Not\_Initialized} when a session is used without
being held or set up; the AES driver's \texttt{Supported\_Key} precondition
(\S\ref{sec:contracts}) rejects a wrong-length key; the range and predicate
subtypes (\S\ref{sec:subtypes}) raise \texttt{Constraint\_Error} on a bad value
under \texttt{-gnata}; and the ext4 filesystem reports I/O and corruption errors
by raising. The lock-free drivers instead report failure through a \texttt{Status}
enum or an \texttt{out Boolean}, precisely so they work on \emph{every} profile,
which is the next point.

\textbf{The cost (and even the availability) depends on the profile.} This
is the embedded-specific part, and it differs across all three:

\begin{itemize}
  \item \texttt{light-tasking} is \texttt{No\_Exception\_Propagation}: there is no
        unwinder, so a \texttt{raise} not handled \emph{in the same subprogram}
        cannot propagate. In practice it reaches the last-chance handler and
        resets the board. The range and predicate checks still \emph{detect}
        errors; you simply cannot catch and recover across calls. Design for
        ``detect, log, reset,'' and report an expected failure with a status
        result, not an exception.
  \item \texttt{embedded} adds \emph{zero-cost} exceptions (ZCX): full propagation,
        handlers, \texttt{Ada.Exceptions} names and messages, and finalization on
        the way out (so an RAII handle releases during an unwind). A handler that
        is not entered costs nothing on the happy path: the cost lives in static
        tables, not in instructions. This is the minimum profile for the RAII
        drivers and the filesystem; here you catch and recover normally.
  \item \texttt{full} has the \emph{same} exception model as \texttt{embedded},
        plus the complete tasking runtime, so exceptions interoperate with dynamic
        tasks, rendezvous and \texttt{abort}, and a stack overflow caught by the
        runtime's guard can be delivered as a catchable \texttt{Storage\_Error}
        instead of a reset (see the Debugging chapter).
\end{itemize}

When an exception escapes everything (always on light-tasking, and as a final
backstop elsewhere) the runtime's \emph{last-chance handler} runs; on this port
it reports the fault and resets, and the \texttt{esp32s3\_safe\_state} example
shows a fault being classified there (see the Debugging chapter).

\textbf{Strategy on a microcontroller.} Prefer the cheap, local mechanisms to
\emph{prevent} bad values: range and predicate subtypes (\S\ref{sec:subtypes}),
\texttt{'Valid} (\S\ref{sec:valid}), contracts (\S\ref{sec:contracts}); use a
status result for an expected, recoverable condition (a NACK, a CRC mismatch); and
reserve exceptions for the genuinely exceptional, remembering that on the leanest
profile ``handling'' one means a controlled reset, not a resumption.

\textbf{A runnable demonstration.} The \texttt{esp32s3\_exceptions} example
exercises all four outcomes on the chip: local handling, propagation across a
call, a re-raise to an outer handler, and an unhandled exception reaching the
last-chance handler (it installs a custom one, exported as
\texttt{\_\_gnat\_last\_chance\_handler}, that prints the exception and halts).

\section{Restricting the language: \texttt{pragma Restrictions} and profiles}
\label{sec:restrictions}\index{Restrictions pragma}

The flip side of choosing features is forbidding them. \texttt{pragma
Restrictions} lets you ban constructs that have no place in tiny or certifiable
firmware and have the compiler \emph{prove} you stayed within the subset:

\begin{lstlisting}
   pragma Restrictions (No_Dependence => Ada.Unchecked_Deallocation);
   pragma Restrictions (No_Recursion, No_Secondary_Stack);
   pragma Restrictions (No_Exceptions);          -- if a raise would be fatal anyway
\end{lstlisting}

A named bundle of these is a \emph{profile}: \texttt{pragma Profile (Ravenscar)}
(and its superset \texttt{Jorvik}) selects the analysable tasking subset
(no dynamic task creation, no entry queues longer than one, a single delay form),
which is what the \texttt{light-tasking} and \texttt{embedded} runtimes are built
against. The Runtime Profiles chapter covers the runtime side; the point here is
that the restriction is a \emph{language} statement the compiler enforces unit by
unit, so a violation is a build error, not a surprise on silicon.

\section{Controlling elaboration: \texttt{Pure}, \texttt{Preelaborate}, \texttt{No\_Elaboration\_Code}}
\label{sec:elab}\index{elaboration}

Elaboration (\S\ref{sec:terms}) is the code that runs before \texttt{Main} to
bring declarations to life. On a microcontroller you usually want \emph{less} of
it: less start-up code, and data that can live in flash rather than be built in
RAM at boot. Category pragmas let you state (and have the compiler enforce)
how little a unit needs:

\begin{itemize}
  \item \texttt{pragma Pure:} the unit has no state and no elaboration at all
        (only pure functions and constants); the strongest promise, and its data
        can sit in flash.
  \item \texttt{pragma Preelaborate:} the unit may be elaborated without
        running any code: all its data has static initial values, so there is no
        boot-time cost and nothing is built at start-up.
  \item \texttt{pragma Elaborate\_Body:} forces the body to be elaborated
        immediately after the spec, a sequencing guarantee that keeps elaboration
        order predictable.
  \item \texttt{pragma Restrictions (No\_Elaboration\_Code):} the compiler
        \emph{rejects} any unit that would emit elaboration code: a hard guarantee
        that start-up is free.
\end{itemize}

\begin{lstlisting}
package Sensor_Tables is
   pragma Pure;                          -- pure: a flash-resident lookup table
   Lut : constant array (0 .. 255) of Integer := ( ... );
end Sensor_Tables;

package Ring is
   pragma Preelaborate;                  -- no code runs to elaborate this unit
   ...
end Ring;
\end{lstlisting}

\textbf{Embedded payoff.} A \texttt{Pure} or \texttt{Preelaborate} unit
contributes nothing to boot time and keeps its constants out of RAM; the runtime
itself builds its lowest layers with \texttt{No\_Elaboration\_Code} so the kernel
has no hidden start-up. The pragmas turn ``this should be cheap to start'' into a
checked property.

